% !TEX root = ../main.tex
\section{Related Work}
\label{sec:related_work}

\subsection{Agent Skills and Skill Benchmarks}
\label{sec:rw_skills}

Skills extend language-model agents with reusable procedural knowledge: task descriptions, invocation conditions, parameter schemas, and multi-step workflows.
\citet{skillsbench2026} introduce \textsc{SkillsBench}, evaluating 86~tasks across 11~domains.
Curated skills boost agent performance by an average of 16.2~percentage points, with the medical domain showing the largest gain (+51.9~pp), while self-generated skills yield near-zero improvement.
\citet{skillnet2026} scale the effort to over 200K~skills with \textsc{SkillNet}, proposing multi-dimensional \emph{static} evaluation along Safety, Completeness, Executability, Maintainability, and Cost-awareness.

Both contributions establish that skill quality matters.
However, neither grounds its evaluation in \emph{execution traces}: \textsc{SkillsBench} measures task-level delta, not whether the skill was actually invoked and followed; \textsc{SkillNet} scores interface quality without running the skill in an isolated sandbox.
\medskillbench bridges this gap by evaluating 1,462~medical skills (917~Tier-A) at the execution level, quantifying the \rbu gap between reading a \skillmd and correctly executing it.

\subsection{Medical and Biomedical Agent Benchmarks}
\label{sec:rw_medical_bench}

Recent medical benchmarks have diversified the evaluation surface for language-model agents.
\citet{medagentbench2025} present \textsc{MedAgentBench}, comprising 300~FHIR-compliant EHR tasks that test interactive clinical actions.
\citet{healthbench2025} release \textsc{HealthBench} with 5,000~multi-turn health dialogues scored by 262~physician-authored rubrics.
In biomedical research, \citet{labbench2024} propose \textsc{LAB-Bench}, covering 2,400+~biology research questions spanning LitQA, DBQA, and ProtocolQA.
\citet{hle2025} introduce the \textsc{Humanity's Last Exam} (HLE), a 2,500-question expert-level benchmark published in \emph{Nature}, whose Biomedicine subset---together with \textsc{LAB-Bench}---has proven effective for evaluating biomedical agents \citep{stella2025}.
Interactive tool-use benchmarks such as $\tau$-bench \citep{taubench2024} and \textsc{ToolSandbox} \citep{toolsandbox2024} further formalize stateful, multi-turn evaluation with programmatic verification.
\citet{medagentgym2025} contribute \textsc{MedAgentGym}, framing medical agent training as a reinforcement learning problem over clinical environments.

These benchmarks assess \emph{task-level} success.
\medskillbench operates at a complementary \emph{skill level}: given a specific \skillmd, does the agent read, invoke, and correctly follow the skill specification?

\subsection{Tool and Skill Description Optimization}
\label{sec:rw_tool_opt}

A growing body of work shows that how tools are \emph{described} is itself a performance lever.
\citet{easytool2024} compress verbose, heterogeneous tool documentation into concise, standardized instructions (\textsc{EasyTool}).
\citet{play2prompt2025} let agents interactively explore tools to auto-generate high-quality usage demonstrations (\textsc{Play2Prompt}).
\citet{tracefree2026} address settings where execution traces are unavailable, using curriculum transfer from trace-rich to trace-free environments to generalize tool-description rewriting (\textsc{Trace-Free+}).
The ReAct paradigm \citep{react2023} further underscores the tight coupling between reasoning traces and tool invocation fidelity.

These methods target \emph{generic} tool descriptions---typically API signatures or short docstrings.
Medical skills are substantially more complex: they encode domain terminology, parameter schemas, safety boundaries, dependency constraints, and multi-step clinical workflows within a single \skillmd.
Our benchmark provides the diagnostic signal needed to evaluate such optimization methods on full-length, structured medical skill documents.

\paragraph{Positioning.}
\Cref{tab:rw_comparison} summarizes how \medskillbench relates to prior work across five key dimensions.

\begin{table}[t]
\centering
\small
\setlength{\tabcolsep}{3pt}
\caption{Comparison of \medskillbench with related work. \emph{Eval.}: task-level (T) or skill-level (S); parenthetical indicates execution-based (E) or static (S) evaluation. \emph{Opt.\ Target}: what is optimized. \emph{Scale}: number of evaluation items.}
\label{tab:rw_comparison}
\begin{tabular}{@{}lcccr@{}}
\toprule
\textbf{Work} & \textbf{Eval.} & \textbf{Domain} & \textbf{Opt.\ Target} & \textbf{Scale} \\
\midrule
SkillsBench     & T (S)  & General   & ---            & 86 \\
SkillNet        & S      & General   & ---            & 200K+ \\
MedAgentBench   & T (E)  & Clinical  & ---            & 300 \\
HealthBench     & T (E)  & Health    & ---            & 5,000 \\
LAB-Bench       & T (S)  & Biomed.   & ---            & 2,400+ \\
HLE             & T (S)  & Multi     & ---            & 2,500 \\
EasyTool        & ---    & General   & Tool desc.     & --- \\
Play2Prompt     & ---    & General   & Tool desc.     & --- \\
Trace-Free+     & ---    & General   & Tool desc.     & --- \\
\midrule
\textbf{\medskillbench} & \textbf{S (E)} & \textbf{Medical} & \textbf{---} & \textbf{1,462} \\
\bottomrule
\end{tabular}
\end{table}
